\chapter{Conclusion and Future Work}
This concluding chapter summarizes the overarching achievements and technical contributions presented throughout this thesis. It critically reflects on the practical constraints and structural limitations encountered during the implementation phase and outlines potential avenues for future research to further refine the management of flexible-grid optical networks within TeraFlowSDN.
\section{Summary of Contributions}
This thesis presented a comprehensive integration of a parallel-link-supported Routing and Spectrum Allocation (RSA) framework within the TeraFlowSDN (TFS) ecosystem. The core objective of managing parallel optical links and enhancing physical layer visibility was achieved through targeted architectural, algorithmic, and interface-level developments. The primary contributions of this research are summarized as follows:
\begin{itemize}
    \item \textbf{Enhancement of Device Drivers:} Substantial modifications were implemented within the OpenConfig \texttt{(OCDriver)} and OpenROADM drivers to extract foundational device metadata, including hardware versions, software versions, and manufacturer nomenclature. The parsing mechanisms were fundamentally overhauled using a two-phase extraction process and model-based filtering, allowing the controller to dynamically compute and store critical RSA parameters—such as minimum/maximum frequencies and flex grid slots—directly into an extended Context Service database schema.
    \item \textbf{ParallelOpticalController and RSA Implementation:}\\ \texttt{ParallelOpticalController} microservice was developed and deployed as a core component of the TFS ecosystem to manage complex optical intents. By leveraging NetworkX to construct multi-directed graphs, the architecture successfully modeled multi-edge parallel optical links. The path computation engine utilized a stateful, bitmap-based intersection methodology adhering to the ITU-T G.694.1 6.25 GHz slot granularity to accurately discover, compute, and allocate contiguous spectrum spaces across varying topologies.
    \item \textbf{UI and UX Enhancements:} The default TFS WebUI was significantly augmented to provide comprehensive spectrum visualization. A dynamic Spectrum Allocation Table was introduced to display explicit slot availability (categorized as \texttt{AVAILABLE}, \texttt{UNAVAILABLE} and \texttt{IN\_USE}) derived directly from endpoint frequency capabilities. Furthermore, parallel links were visually abstracted using a numeric notation within the topology viewer, ensuring clear topological representation without graphical overflow.
\end{itemize}
Collectively, these structural and algorithmic enhancements facilitated a highly responsive and deterministic optical control plane. The computational performance and spectrum utilization efficiency evaluations demonstrated that the bitmap-based allocation algorithm efficiently handles heavy topological data while minimizing spectrum fragmentation. Furthermore, the throughput analysis—validated through rigorous testing of network-level traffic loads—confirmed that the system maintains an optimal blocking probability by dynamically distributing slot allocations across all available parallel links. While the proposed solutions achieved strict functional validation and showcased strong computational efficiency, the development and physical testbed deployment phases surfaced specific architectural and environmental constraints. These systemic boundaries define the operational envelope of the current controller architecture and introduce several critical limitations that must be addressed.
\section{Limitations}
While the integration of the \texttt{ParallelOpticalController} introduces significant capabilities to the TeraFlowSDN (TFS) ecosystem—particularly in managing parallel optical links and rendering detailed spectrum visualizations—the current framework inherently operates within specific systemic and architectural boundaries. Acknowledging these limitations is critical for understanding the operational envelope of the proposed solution and identifying areas for continuous research. The constraints can be broadly categorized into architectural dependencies, algorithmic assumptions, and device-level interactions.
\subsection*{\small Architectural Dependencies and Performance Overheads}
A fundamental architectural limitation is the stateful operation of the \texttt{ParallelOptic-\\alController}. By design, microservices are traditionally engineered to be stateless and independently executable. However, the proposed RSA computation engine is deeply ingrained within the TFS control plane, heavily relying on continuous synchronization with the Context Service to fetch the latest topology, optical configurations, and device capabilities. While statefulness is a common paradigm in centralized Software-Defined Networking (SDN) architectures, it tightly couples the controller to the TFS database schema, limiting its standalone portability. Furthermore, this state dependency introduces significant performance overheads. The current implementation relies on a quantitative volume of database queries to extract granular endpoint indexes, associated channel mappings, and flex-grid bitmap values across the multi-directed graph. This iterative polling of the database during the \texttt{find\_paths()} and \texttt{expand\_path()} execution phases introduces computational latency. In highly scaled networks with dense parallel links, this latency could scale linearly, potentially impacting real-time provisioning SLA requirements.

\subsection*{\small Algorithmic and Topological Assumptions}
The algorithmic approach to graph modeling and path discovery introduces several operational constraints. First, the \texttt{NetworkX} graph construction \texttt{(G\_FREE and G\_simple\_free)} fundamentally assumes that all optical links are bi-directional. The Dijkstra-based shortest path and simple path expansion algorithms compute routing metrics based on this symmetry. If the physical infrastructure contains uni-directional optical fibers or asymmetric spectrum availability, the current algorithmic formulation will fail to discover a valid end-to-end lightpath. Second, there is a topological abstraction limitation regarding ROADM internal architectures. The graph expansion mechanism assumes that any parallel links connecting a Transponder to a ROADM, or between two ROADMs, are functionally equivalent as long as they reach the target node. However, it currently does not account for the specific spatial Degrees (DEGs) to which these fibers terminate. As observed in the roadm.xml definitions (e.g., ports designated to D1 versus D2), if a fiber terminates on a different DEG, the internal optical cross-connect (OXC) constraints of the ROADM might prevent successful switching. The controller currently abstracts this away, treating any parallel link to the same device as a valid routing option, which could lead to physically unrealizable path computations.

\subsection*{\small Lifecycle Management and Device-Level Interactions}
From an operational lifecycle perspective, the current framework focuses primarily on the establishment of optical intents. The slot acquisition process currently requires manual intervention and intuitive oversight by the network operator. More importantly, the controller does not natively support an automated lightpath teardown and slot release mechanism. Once slots are marked as \texttt{IN\_USE} in the spectrum bitmap, reversing this hash to free the contiguous 6.25 GHz slots upon intent deletion requires manual reconciliation. Finally, the device interaction layer remains bounded by the scope of the testbed. The enhanced information extraction mechanisms within OCDriver and OpenROADMDriver (utilizing filtering modules and targeted model-based extraction) were engineered and validated against specific emulated environments (e.g., SSSA-CNIT Mellanox emulators and virtual ROADMs). This means the extraction logic is heavily tailored to these configurations. A wider array of device configurations and vendor-specific YANG models must be studied to provide generalized, built-in support for diverse physical hardware. Because the controller has not yet been validated against real, physical transponders and ROADMs in a hardware testbed, unforeseen hardware idiosyncrasies, netconf session timeouts, or proprietary XML formatting could introduce unknown errors. Additionally, while the RSA algorithm successfully computes the path and updates the logical SDN database, it lacks the southbound provisioning logic to actively push the required cross-connect configurations down to the physical devices' local databases. Currently, the solution remains fully analytical and logical within the SDN control plane.
\section{Future Work Directions}
Building upon the established framework and the identified limitations, several highly promising avenues for future research and development emerge. The transition from a logical, stateful microservice to an autonomous, hardware-validated optical controller requires focusing on the following key areas:
\subsection*{\small Vendor-Agnostic Device Abstraction and Hardware Validation}
Transitioning the current framework from emulated environments to physical hardware testbeds is a paramount next step. Future iterations of this research should expand the OCDriver and OpenROADMDriver to encompass a broader array of commercial vendor equipment (e.g., Cisco, Nokia, Ciena) and their respective, complex YANG data models. By building a generalized, vendor-agnostic device abstraction layer, the controller can seamlessly adapt to diverse hardware idiosyncrasies. Furthermore, the framework must be extended to include active southbound provisioning. Rather than merely updating the logical state within the TeraFlowSDN database, future work must focus on actively pushing the computed optical cross-connects and frequency configurations directly to the physical network elements via NETCONF.
\subsection*{\small Algorithmic Refinement for Spatial and Directional Constraints}
To accurately reflect the physical realities of modern transport networks, the RSA algorithm must be decoupled from its assumption of strict bi-directionality. Future work should enhance the \texttt{NetworkX} graph modeling to natively support uni-directional optical fibers and asymmetric spectrum bands. Crucially, the path discovery and expansion algorithms must be upgraded to comprehend ROADM internal architectures. Incorporating node-internal spatial constraints—specifically validating that parallel fibers terminating on different spatial Degrees (DEGs) can establish a valid cross-connect—will ensure that computed lightpaths are physically viable at the photonics level.
\subsection*{\small Automated Lifecycle and Spectrum Management}
To achieve a fully autonomous, intent-based network ecosystem, the controller must evolve to manage the complete lifecycle of an optical intent. Research must be directed toward fully automating the slot acquisition process, removing the need for manual operator intervention. More importantly, dynamic lightpath teardown and slot release mechanisms must be implemented. Developing algorithms to intelligently reverse the spectrum bitmap hashes, automatically freeing and coalescing contiguous 6.25 GHz slots upon intent deletion, will be critical for preventing spectrum fragmentation over time.
\subsection*{\small Architectural Optimization for Low-Latency Computation}
To mitigate the latency introduced by the high volume of sequential database queries during the graph expansion and RSA computation phases, the controller's architecture requires optimization. Future research should investigate the integration of an in-memory, distributed caching layer (such as Redis) directly accessible by the ParallelOpticalController. By caching the topological state, device capabilities, and current spectrum bitmaps, the computation engine can reduce its strict dependence on the centralized Context Service. This architectural shift would drive the microservice closer to a stateless operational paradigm, significantly reducing path computation times and improving overall scalability in dense, highly parallelized network topologies.

In conclusion, this thesis successfully bridges a critical gap in the TeraFlowSDN ecosystem by introducing robust support for parallel optical links and granular spectrum management. Through the development of the Parallel Optical Controller microservice and the targeted enhancement of device-level data extraction, the research demonstrates a highly capable approach to Routing and Spectrum Allocation (RSA) within a modern software-defined architecture. While specific architectural dependencies, algorithmic assumptions, and the need for physical hardware validation delineate the boundaries of the current framework, the foundation laid by this work provides a vital stepping stone for continuous development. Ultimately, these structural and algorithmic contributions significantly advance the capacity of SDN controllers to visualize, compute, and orchestrate complex, next-generation optical transport networks with enhanced precision and reliability.